\documentclass[11pt]{article}

\usepackage[T1]{fontenc}
\usepackage[utf8]{inputenc}
\usepackage{mathptmx}
\usepackage{amsmath,amssymb}
\usepackage[a4paper,margin=1in]{geometry}
\usepackage{xcolor}
\usepackage{setspace}
\usepackage{enumitem}
\usepackage{fancyhdr}
\usepackage[hidelinks]{hyperref}

\definecolor{responseblue}{RGB}{0,70,140}
\definecolor{changeblue}{RGB}{0,90,170}
\definecolor{commentgray}{RGB}{70,70,70}

\setstretch{1.08}
\setlength{\parindent}{0pt}
\setlength{\parskip}{0.55em}
\setlength{\headheight}{14pt}

\pagestyle{fancy}
\fancyhf{}
\lhead{Paper-TW-Feb-26-0360}
\rhead{Response to the Editor and Reviewers}
\cfoot{\thepage}

\newcounter{reviewcomment}[section]
\renewcommand{\thereviewcomment}{\arabic{reviewcomment}}

\newenvironment{quotedcomment}
{%
  \begin{list}{}{
    \setlength{\leftmargin}{1.5em}%
    \setlength{\rightmargin}{0.5em}%
    \setlength{\topsep}{0.25em}%
    \setlength{\parsep}{0.35em}%
  }
  \item\relax\color{commentgray}\itshape
}
{\end{list}}

\newcommand{\commentheading}[1]{%
  \refstepcounter{reviewcomment}%
  \par\medskip\noindent\textbf{Comment \thereviewcomment\ (#1)}
}

\newcommand{\generalheading}[1]{%
  \par\medskip\noindent\textbf{#1}
}

\newenvironment{authorresponse}
{\par\noindent\textcolor{responseblue}{\textbf{Response from the Authors:}}\par\color{responseblue}}
{\par}

\newenvironment{manuscriptchanges}
{\par\noindent\textcolor{changeblue}{\textbf{Changes in the Revised Manuscript:}}\par\color{changeblue}}
{\par}

\newcommand{\revisionlocation}[1]{%
  \par\noindent\textbf{Location in the Revised Manuscript:} #1\par
}

\newcommand{\blankreply}{%
  \begin{authorresponse}
  % To be completed after the corresponding revision is finalized.
  \vspace{1.2\baselineskip}
  \end{authorresponse}
  \begin{manuscriptchanges}
  % Insert the exact revised text, equation, figure, table, or a concise description here.
  \vspace{1.2\baselineskip}
  \end{manuscriptchanges}
  \revisionlocation{}
}

\begin{document}

\begin{center}
  {\Large\bfseries Responses to the Comments of the Editor and Reviewers}\par
  \vspace{0.8em}
  {\large\bfseries MEET-COBRA: Machine-learning-based Energy-Efficient proacTive\\
  Coordination of handOver, Beamforming and Resource Allocation}\par
  \vspace{0.6em}
  {\large Manuscript ID: Paper-TW-Feb-26-0360}\par
  \vspace{0.8em}
  Bowen Xie, Sheng Zhou, and Zhisheng Niu\par
  Beijing National Research Center for Information Science and Technology,\par
  Department of Electronic Engineering, Tsinghua University, Beijing 100084, China\par
  E-mail: xbw22@mails.tsinghua.edu.cn, \{sheng.zhou, niuzhs\}@tsinghua.edu.cn
\end{center}

\vspace{0.8em}

Dear Editor and Reviewers,

% Replace this paragraph with the finalized cover note after all revisions are complete.
We sincerely thank the Editor and the Reviewers for their careful evaluation of our manuscript
and for their constructive comments. A concise summary of the major revisions will be provided
below, followed by our detailed point-by-point responses. In the marked-up manuscript, all
changes made in this revision will be highlighted in blue.

\section*{Summary of Major Revisions}

% To be completed after the technical revisions and additional experiments are finalized.
\vspace{4\baselineskip}

\clearpage
\section{Editor}

\setcounter{reviewcomment}{0}
\commentheading{Major revision decision and overall concerns}
\begin{quotedcomment}
While the work is regarded positively in general, the reviews have identified several areas
that need major improvement, e.g., regarding problem formulation, technical rigor/consistency,
evaluation, and presentation. Therefore, I invite you to respond to the reviewer(s)' comments
and revise your manuscript.

It should be emphasized that a major revision decision does not imply eventual acceptance.
A future version of the paper will only be accepted if the revision addresses the major
concerns, and the reply is highly persuasive to the Editors and Reviewers.
\end{quotedcomment}
\blankreply

\clearpage
\section{Reviewer 1}

\setcounter{reviewcomment}{0}
\commentheading{Motivation and cross-tier heterogeneity}
\begin{quotedcomment}
The paper explains the coupling among handover, beamforming, and resource allocation in
HetNets. It would be beneficial to further strengthen the motivation by more clearly
articulating why this problem becomes fundamentally more unique and challenging in
heterogeneous networks, beyond differences in system architecture. In particular, additional
discussion on the new optimization challenges introduced by cross-tier heterogeneity could
further improve the clarity of the problem formulation.
\end{quotedcomment}
\begin{authorresponse}
We agree that the earlier motivation described the architectural differences between the two
tiers without making their optimization consequences sufficiently explicit. We have therefore
expanded the Introduction to explain why a cross-tier HO changes several coupled quantities at
once: carrier and RB configuration, power cost, beam-training requirement, link gain, RB demand,
and interference footprint. We also explain how the frame-level association decision feeds into
slot-level BF and RA through tier-dependent service rates and capacity constraints. This addition
clarifies the specific difficulty of COBRA in a HetNet rather than treating heterogeneity merely
as a deployment detail.
\end{authorresponse}
\begin{manuscriptchanges}
We added the following paragraph:

\emph{``The cross-tier heterogeneity creates an optimization problem that is qualitatively
different from association among homogeneous BSs. A macro link offers robust wide-area coverage
but has a different carrier frequency, RB bandwidth, power cost, and antenna configuration from
a directional mmWave micro link. A handover therefore changes not only the serving BS but also
the available RB budget, beam-training requirement, desired-link gain, and interference footprint.
Moreover, directing a micro-BS beam toward one vehicle changes the RB demand of that link, while
the resulting load and RB occupancy affect the interference experienced by vehicles in other
micro cells. Consequently, frame-level cross-tier association and slot-level BF/RA decisions are
coupled through tier-dependent service rates and capacity constraints.''}
\end{manuscriptchanges}
\revisionlocation{Pages 1--2, Section I (Introduction).}

\commentheading{Queueing delay and end-to-end latency}
\begin{quotedcomment}
The latency formulation is currently based on queue-length constraints. The paper could be
further strengthened by discussing the relationship between queueing delay and actual
end-to-end latency in a two-tier vehicular HetNet. Incorporating or discussing additional
latency components, such as handover interruption, beam training overhead, retransmissions,
and signaling delay, would improve the realism of the latency evaluation.
\end{quotedcomment}
\begin{authorresponse}
We agree that the original presentation could be read as equating an instantaneous scaled
backlog with actual end-to-end packet latency. That interpretation is not justified by Little's
law. We have corrected the formulation and terminology throughout the manuscript. The revised
paper now states that Little's law is a long-term average relation and defines
$d_v^{\rm Q}(n)=q_v(n)/\lambda_v$ only as a normalized-backlog queueing-delay proxy. Accordingly,
$U$ is now described as a queue-length constraint violation probability, not as a hard per-packet
or end-to-end delay guarantee.

We have also delineated the modeled latency components. Beam-search overhead affects the service
rate through $\eta_{m,v}^{[x]}(i)$, whereas HO interruption, control signaling, processing,
packetization, and retransmission delays are not included. We explicitly state this scope instead
of implying that the present metric establishes uRLLC-grade end-to-end reliability. The abstract,
problem statement, result discussion, figure captions, and the on-figure axis labels have been
revised consistently.
\end{authorresponse}
\begin{manuscriptchanges}
The revised definition and scope statement read:

\emph{``Little's law relates the long-term average backlog and average queueing delay under
stationarity; it does not provide a sample-path identity between the instantaneous backlog and
the delay of an individual packet. Motivated by this average relation, we define the
normalized-backlog quantity $d_v^{\rm Q}(n)\triangleq q_v(n)/\lambda_v$ as a queueing-delay
proxy.''}

\emph{``Accordingly, $U$ measures normalized-backlog violations rather than a hard per-packet or
end-to-end delay guarantee. Beam-search overhead is included in the service rate through
$\eta_{m,v}^{[x]}(i)$; however, HO interruption, control signaling, processing, packetization,
and retransmission delays are outside the present model.''}
\end{manuscriptchanges}
\revisionlocation{Page 5, Section II-C, Eqs. (22)--(24); Pages 12--13, Figs. 6--7 and Section V-D; Page 14, Section VI.}

\commentheading{Terminology in Fig. 1}
\begin{quotedcomment}
In Fig. 1, the use of the term ``Optimal'' beam-pair and gain prediction may be slightly
strong given that the neural network provides estimated outputs. The presentation could be
made more precise by adopting terminology such as ``Predicted'' or ``Estimated'' beam-pair/gain.
\end{quotedcomment}
\begin{authorresponse}
We agree and have removed the potentially misleading use of ``Optimal'' for an NN output. Fig. 1
now identifies the outputs as next-frame beam-pair probabilities and channel-gain predictions.
The accompanying text distinguishes the predicted probability distribution from the
ground-truth best-beam label used during offline training. For consistency, the corresponding
terminology in Fig. 4 and its discussion has also been changed to ``best-beam desired-link gain.''
\end{authorresponse}
\begin{manuscriptchanges}
The revised Fig. 1 caption is:

\emph{``The architecture of the proposed neural network models for next-frame beam-pair
probability and channel-gain prediction.''}

The revised text states that the model predicts the \emph{``next-frame beam-pair probability
distribution, best-beam desired-link gain, and interference-link gain.''}
\end{manuscriptchanges}
\revisionlocation{Page 6, Fig. 1 and Section IV-A; Page 11, Fig. 4 and Section V-B.}

\commentheading{Propagation of prediction errors}
\begin{quotedcomment}
The paper presents promising prediction accuracy results. It would be valuable to further
analyze how prediction errors propagate to higher-layer decisions, particularly proactive
handover performance. Additional experiments quantifying the impact of prediction errors on
HO reliability, association accuracy, or latency performance would provide deeper insight
into the robustness of the proposed framework.
\end{quotedcomment}
\blankreply

\commentheading{Notation table}
\begin{quotedcomment}
The notation is relatively dense, and multiple related symbols are introduced throughout the
paper. Adding a notation table summarizing the main variables and symbols would significantly
improve readability and make the paper easier to follow.
\end{quotedcomment}
\begin{authorresponse}
We agree that the original notation density made the model unnecessarily difficult to follow.
We have added a two-column notation table that groups the principal sets and indices, time and
association variables, channel/BF quantities, queue and resource variables, prediction outputs,
and algorithm parameters. Symbols used only locally remain defined next to their equations so
that the table remains compact.
\end{authorresponse}
\begin{manuscriptchanges}
We added Table I, entitled \emph{``Summary of Main Notation.''}
\end{manuscriptchanges}
\revisionlocation{Page 4, Table I, at the beginning of Section II.}

\commentheading{Selection and sensitivity of heuristic parameters}
\begin{quotedcomment}
Several heuristic parameters, such as $\zeta$, $M_{\rm P}$, and $N_{\rm iter}$, are introduced
in the framework. Providing additional discussion on their selection criteria, practical
interpretation, or parameter sensitivity would further improve the reproducibility and
robustness evaluation of the proposed method.
\end{quotedcomment}
\begin{authorresponse}
We have added both the numerical settings and their operational interpretations. In particular,
$M_{\rm P}=5$ is the point at which the measured top-$M_{\rm P}$ beam reliability already exceeds
99\%, while every additional candidate incurs another probe; $\zeta=1.1$ relaxes an infeasible RB
budget in moderate 10\% increments; $\epsilon=1$ RB matches the granularity of the final integer
allocation; and $N_{\rm iter}=3$ bounds the frame-level fixed-point computation. These values are
now also listed in the simulation-parameter table, which makes the setup directly reproducible.
\end{authorresponse}
\begin{manuscriptchanges}
We added the following explanation:

\emph{``The heuristic parameters have direct algorithmic interpretations. We use
$M_{\rm P}=5$ because the measured top-$M_{\rm P}$ beam reliability already exceeds 99\% around
this operating point, while each additional candidate incurs another beam probe. The relaxation
factor $\zeta=1.1$ expands an infeasible RB budget in moderate 10\% steps rather than making a
single aggressive relaxation. We set $\epsilon=1$ RB, matching the granularity of the final
integer allocation, and cap the fixed-point refinement at $N_{\rm iter}=3$ to bound the
frame-level control computation.''}
\end{manuscriptchanges}
\revisionlocation{Page 8, Section IV-B; Page 11, Table II.}

\commentheading{Neural-network complexity and inference overhead}
\begin{quotedcomment}
The paper emphasizes the lightweight nature of the ML inference process. Including
quantitative information regarding neural network complexity, such as model size, inference
latency, FLOPs, or runtime overhead, would further strengthen the practicality claims of the
proposed framework.
\end{quotedcomment}
\blankreply

\clearpage
\section{Reviewer 2}

\setcounter{reviewcomment}{0}
\generalheading{Overall Assessment and Strengths}
\begin{quotedcomment}
The joint treatment of proactive handover, beamforming, and resource allocation is a real and
underexplored gap, and using the per-link resource-block demand as the common currency that
links prediction, association, and scheduling is clean and well motivated. The evaluation
platform is credible, since combining SUMO traces with Sionna RT ray tracing is more realistic
than the stochastic channel models most comparable papers use. The algorithmic engineering is
careful, with per-module complexity analysis, a Shmoys--Tardos approximation for the assignment
step with a stated two-factor capacity bound, and a reasonable ablation and oracle design.
Theorem 1 gives a principled starting point for the demand estimate. Overall the work is
competent and the writing is mostly clear.
\end{quotedcomment}
\begin{authorresponse}
We thank the Reviewer for recognizing the value of the three-way coordination framework, the
RB-demand coupling, the SUMO/Sionna RT evaluation platform, and the algorithmic analysis. We
have retained these elements and revised the manuscript carefully to address the formulation,
evaluation, and presentation concerns raised below.
\end{authorresponse}
\revisionlocation{No manuscript change is required for this overall assessment.}

\commentheading{Interference-model consistency and validation}
\begin{quotedcomment}
The interference model is crude and inconsistent with the proposed algorithms. The
interference-link gain in Eq. (10) uses a Frobenius-norm average, justified by random transmit
and receive beam directions, and the resource-block overlap term in Eq. (16) assumes
independent and uniformly random allocation across base stations. Neither holds in the
proposed system, because every interfering micro base station runs the same deterministic
PET-BF, directing beams at its own vehicles, and the same deterministic OTR-RA, which places
resource blocks by efficiency rather than at random. Since the SINR in Eq. (17), and therefore
the entire power and latency story, rests on this approximation, please quantify how far these
surrogates deviate from directional interference, for example against a beam-aware interference
computation from Sionna RT, or justify why the bias is negligible at the reported loads.
\end{quotedcomment}
\blankreply

\commentheading{Queue proxy versus delay guarantee}
\begin{quotedcomment}
The latency constraint is a mean-rate queue proxy, not a delay guarantee, yet the results are
framed as uRLLC-grade reliability. Section II-C invokes Little's law to write the per-slot
constraint in Eq. (23). Little's law is a long-run average identity, not a per-slot map from
instantaneous backlog to instantaneous delay, so Eq. (23) bounds a scaled queue length rather
than actual packet latency, and the metric $U$ in Eq. (24) inherits this. Given the repeated
uRLLC framing in the introduction, please either justify this rigorously.
\end{quotedcomment}
\begin{authorresponse}
We agree with the Reviewer's technical point. Little's law cannot justify a per-slot identity
between $q_v(n)/\lambda_v$ and the delay of an individual packet, and we do not attempt to defend
that interpretation in the revision. Instead, we have corrected the model statement: Little's
law is used only as motivation for a normalized-backlog proxy
$d_v^{\rm Q}(n)=q_v(n)/\lambda_v$, and Eq. (23) is explicitly described as a queue-length-based
QoS constraint. Consequently, $U$ in Eq. (24) is a queue-length constraint violation probability,
not a delay-reliability guarantee.

We have removed the uRLLC-grade/end-to-end interpretation and narrowed the Introduction to a
single class of delay-sensitive downlink traffic. We also list the latency components excluded
from the present model and identify the one overhead component already included: beam-search
overhead enters the service rate through $\eta_{m,v}^{[x]}(i)$. The abstract, formulation,
result text, captions, and on-figure axis labels have all been made consistent with this corrected
interpretation. Thus, the numerical results themselves are unchanged, but the claim attached to
them is now the claim that the simulated metric can support.
\end{authorresponse}
\begin{manuscriptchanges}
We replaced the original instantaneous-delay interpretation with:

\emph{``Little's law relates the long-term average backlog and average queueing delay under
stationarity; it does not provide a sample-path identity between the instantaneous backlog and
the delay of an individual packet. Motivated by this average relation, we define the
normalized-backlog quantity $d_v^{\rm Q}(n)\triangleq q_v(n)/\lambda_v$ as a queueing-delay
proxy.''}

We further clarify:

\emph{``Accordingly, $U$ measures normalized-backlog violations rather than a hard per-packet or
end-to-end delay guarantee. $\ldots$ The reported $d_v^{\rm Q}$ and $U$ should therefore be
interpreted as radio-access queueing proxies, not as uRLLC-grade end-to-end reliability.''}
\end{manuscriptchanges}
\revisionlocation{Page 1, Abstract and Section I; Page 5, Section II-C and Eqs. (23)--(24); Pages 12--13, Figs. 6--7 and Section V-D; Page 14, Section VI.}

\commentheading{Impact of gain-prediction error}
\begin{quotedcomment}
Gain-prediction error is not propagated into the performance metrics, and 4 dB MAE is large
for mmWave link budgeting. Figure 4(b) reports validation MAE around 4.07 dB for the
desired-link gain. At mmWave, 4 dB can move a link across the boundary between a usable and
unusable beam, and it perturbs the demand estimate in Eq. (38) and hence $U$. Please add an
error-propagation study, or at least plot $U$ against injected gain-prediction noise, to show
the accuracy achieved is actually sufficient.
\end{quotedcomment}
\blankreply

\commentheading{Generalization of the learned models}
\begin{quotedcomment}
Generalization of the learned models is untested. The network is trained and validated on
6000 frames from a single map, a fixed four-micro-BS layout, and one traffic regime, with the
split taken over non-overlapping time intervals of the same scenario, so all reported accuracy
is in-distribution. There is no evaluation under a different topology, base-station placement,
or arrival pattern, and the limitation is not acknowledged beyond a vague mention of online
adaptation. Please add an out-of-scenario test, or state plainly that the model is site-specific
and must be retrained per deployment.
\end{quotedcomment}
\begin{authorresponse}
We agree that the existing temporal split does not test cross-site generalization. We have taken
the second route suggested by the Reviewer and now state this limitation plainly. Specifically,
the revised paper identifies the model as site-specific, describes the 70/30 split as
non-overlapping temporal segments from the same map, BS placement, and traffic setting, and calls
the resulting accuracy in-distribution performance. We no longer imply that this experiment
demonstrates zero-shot generalization. The paper also states that a substantially different map
or BS layout would require site-specific fine-tuning or retraining from local measurements.
\end{authorresponse}
\begin{manuscriptchanges}
We added the following scope statement:

\emph{``Its parameters depend on the deployment geometry, BS placement, and channel/traffic
distribution represented in the training traces. Consequently, the temporal validation
conducted in this paper measures in-distribution performance within the same deployment and
does not establish zero-shot generalization to a new map or BS layout. Deployment in a
substantially different environment would require site-specific fine-tuning or retraining using
local channel measurements.''}

The dataset description now explicitly states that the 70/30 training/validation samples are
drawn from non-overlapping temporal segments of the same map, four-micro-BS placement, and
traffic-generation setting.
\end{manuscriptchanges}
\revisionlocation{Page 6, Section IV-A; Page 11, Section V-B.}

\commentheading{Learning-based or optimization-based comparator}
\begin{quotedcomment}
The authors cite learning-based proactive handover and reinforcement-learning beam tracking,
yet none appears as a baseline, and there is no conventional optimization comparator. A
32 percent gain over a deliberately weak baseline is not strong evidence of an advance.
Please add at least one learning-based or optimization-based comparator.
\end{quotedcomment}
\blankreply

\commentheading{Inference, control-plane, and handover overhead}
\begin{quotedcomment}
The paper claims lightweight inference overhead but gives no inference latency, FLOPs, or
control-plane bit budget. More importantly, the central selling point of proactive handover is
reduced failures and interruption, yet handover signaling and interruption latency are deferred
to future work and excluded from $U$. Please add a like-for-like overhead comparison and at
least a first-order handover-latency model.
\end{quotedcomment}
\blankreply

\commentheading{Traffic arrival process and heterogeneous QoS}
\begin{quotedcomment}
The traffic and QoS model is thinner than the motivation promises. The arrival process in
Eq. (22) is characterized only by its mean, with no stated distribution, yet for a study whose
headline metric is a violation tail, arrival burstiness dominates the result and $U$ is hard to
interpret without it. In addition, the introduction motivates heterogeneous QoS across eMBB,
uRLLC, and energy, but the model carries a single backlog queue and a single latency class.
Please specify the arrival process and either add a second QoS class or narrow the motivation.
\end{quotedcomment}
\begin{authorresponse}
We agree on both points. First, the revised model now specifies independent Poisson arrivals:
$a_v(n)$ has mean $\lambda_v\Delta t$ and is independent across vehicles and slots. This makes
the traffic assumption and the interpretation of the reported tail metric explicit. Second,
rather than adding a heterogeneous service model that is not supported by the current algorithms
or experiments, we have narrowed the paper's motivation and scope. The revised Introduction and
traffic model state that we consider one class of delay-sensitive downlink traffic with a common
queueing requirement; differentiated packet classes, deadlines, reliability targets, and
end-to-end uRLLC guarantees are outside the present scope.
\end{authorresponse}
\begin{manuscriptchanges}
The traffic model now states:

\emph{``We model $a_v(n)$ as an independent Poisson random variable with mean
$\lambda_v\Delta t$, where $\lambda_v$ is the average data arrival rate of vehicle $v$;
arrivals are independent across vehicles and slots. The model represents one class of
delay-sensitive downlink traffic with a common queueing requirement. Differentiated packet
classes, deadlines, and reliability targets are not considered.''}
\end{manuscriptchanges}
\revisionlocation{Page 1, Section I; Page 5, Section II-C, immediately below Eq. (22).}

\commentheading{OTR-RA acronym in Fig. 2}
\begin{quotedcomment}
Figure 2 labels the resource-allocation block OUR-RA, while the text and Algorithm 3 use
OTR-RA. Please fix the acronym.
\end{quotedcomment}
\begin{authorresponse}
We have corrected the typographical error. The resource-allocation block in Fig. 2 now reads
``Algorithm 3 (OTR-RA),'' consistent with the text and Algorithm 3.
\end{authorresponse}
\begin{manuscriptchanges}
In Fig. 2, ``OUR-RA'' was replaced by ``OTR-RA.''
\end{manuscriptchanges}
\revisionlocation{Page 10, Fig. 2.}

\clearpage
\section{Reviewer 3}

\setcounter{reviewcomment}{0}
\generalheading{Summary}
\begin{quotedcomment}
This manuscript proposes MEET-COBRA, a framework for jointly coordinating handover,
beamforming, and resource allocation in a two-tier vehicular network, using machine learning
(LSTM) to predict channel gains and beam pairs ahead of time. The authors argue that solving
this joint problem exactly is computationally impractical and instead propose a two-timescale
approach: predictions drive the handover decision once per frame, while beamforming and
resource allocation adapt every slot based on those predictions and on real-time measurements.
The simulation setup combines SUMO mobility traces with Sionna ray-tracing channel data and
the evaluation includes both oracle bounds and per-module ablations that give a fairly complete
picture of where the reported gains come from. The proposed scheme shows a meaningful
reduction in transmit power and latency violations compared to a reactive baseline. I have a
number of concerns detailed below that should be addressed in a revision.
\end{quotedcomment}
\begin{authorresponse}
We thank the Reviewer for the accurate summary and for recognizing the value of the two-timescale
coordination, the SUMO/Sionna RT platform, and the oracle and ablation studies. In response to the
comments below, we have clarified the relationship to our conference paper, expanded the related
work, corrected the queue-based QoS terminology, and improved the presentation and reproducibility.
The comments that require additional overhead, baseline, and mobility studies are being addressed
separately in the technical revision.
\end{authorresponse}
\revisionlocation{The corresponding manuscript changes are identified in the point-by-point responses below.}

\subsection*{Major Comments}

\commentheading{Relationship to the authors' conference paper}
\begin{quotedcomment}
This work reads as an extension of the authors' own earlier conference paper, which focused
on joint handover and resource allocation for a similar vehicular HetNet setup. The framework
here appears to build on that by adding beamforming as a third coordinated decision, along
with a more detailed channel model and simulation setup. Given this overlap, I would encourage
the authors to state clearly, early in the paper, that this is an extended version of that prior
work, and to briefly summarize what is specifically new here, since it would give readers a
complete picture of what this paper contributes.
\end{quotedcomment}
\begin{authorresponse}
We agree that the relationship and incremental technical contribution should be transparent near
the beginning of the paper. We now explicitly identify AEPHORA as our conference paper and list
three substantive extensions: (i) expanding joint PHO/RA to three-way PHO/BF/RA coordination,
including beam-search overhead and inter-cell interference; (ii) replacing mobility-only
prediction with beam-pair, desired-link-gain, and interference-link-gain predictions and adding
the PET-BF and OTR-RA algorithms; and (iii) evaluating the full design with SUMO/Sionna RT MIMO
channels, oracle benchmarks, and module-level ablations. This statement is placed before the main
contribution list so that readers can assess the relationship immediately.
\end{authorresponse}
\begin{manuscriptchanges}
We added:

\emph{``This article substantially extends our conference paper AEPHORA. First, it expands the
decision space from joint PHO/RA to the three-way coordination of PHO, BF, and RA and explicitly
incorporates beam-search overhead and inter-cell interference. Second, it replaces mobility-only
prediction with link-level predictions of beam-pair probabilities, desired-link gains, and
interference-link gains, and introduces the new PET-BF and OTR-RA algorithms around a
generalized-assignment-based PHO procedure. Third, it evaluates the complete framework using
SUMO mobility and MIMO channels generated by Sionna RT, together with oracle benchmarks and
module-level ablations.''}
\end{manuscriptchanges}
\revisionlocation{Page 2, Section I (Introduction), immediately before the problem statement and contribution list.}

\commentheading{ML computation and reporting overhead versus CSI overhead}
\begin{quotedcomment}
One of the paper's central motivations is reducing the overhead of channel state acquisition,
which is a reasonable goal. That said, in the proposed design, each vehicle needs to run three
separate ML prediction models locally every frame, one each for beam-pair selection,
desired-link gain, and interference-link gain, with all of it reported to the macro base station
each frame. This introduces its own computational and reporting overhead, and it isn't clear
from the paper how that overhead compares to the CSI overhead the framework is designed to
avoid in the first place. A brief discussion of this would strengthen the paper's claim of
achieving lower overhead with lightweight ML inference.
\end{quotedcomment}
\blankreply

\commentheading{Coordinated handover and beamforming baseline}
\begin{quotedcomment}
The comparison in Section V is against a reactive baseline, and that baseline is fairly weak on
all three fronts at once: handover triggered by signal-strength thresholds, beam selection
through random probing, and resource allocation under a worst-case interference assumption.
The per-module ablation study shown in the paper is a nice complement to this, since it shows
what each of the three proposed modules is individually contributing. It would also help to
compare against a coordinated handover and beamforming baseline from existing literature.
\end{quotedcomment}
\blankreply

\commentheading{Vehicle-speed distribution and mobility sensitivity}
\begin{quotedcomment}
Despite high-mobility vehicular scenarios being the paper's central motivation, it does not
report any information on vehicle speed anywhere in the manuscript. Speed is arguably one of
the main factors driving how quickly channel conditions change from one frame to the next, so
it directly affects how difficult the prediction problem actually is. The authors do state that
training and validation samples are drawn from non-overlapping time intervals of the same
simulated traffic, but the vehicle speed distribution within that dataset is never reported. It
would help to see the speed distribution used, and a result showing how performance of the
proposed framework changes across different mobility levels.
\end{quotedcomment}
\blankreply

\subsection*{Minor Comments}

\commentheading{OTR-RA acronym in Fig. 2}
\begin{quotedcomment}
Fig. 2 labels the resource allocation block ``OUR-RA,'' while the surrounding text and
Algorithm 3 consistently use ``OTR-RA.''
\end{quotedcomment}
\begin{authorresponse}
We have corrected the typographical error. The resource-allocation block in Fig. 2 now reads
``Algorithm 3 (OTR-RA),'' consistently with Algorithm 3 and the surrounding text.
\end{authorresponse}
\begin{manuscriptchanges}
In Fig. 2, ``OUR-RA'' was replaced by ``OTR-RA.''
\end{manuscriptchanges}
\revisionlocation{Page 10, Fig. 2.}

\commentheading{Handover-algorithm parameters in the simulation table}
\begin{quotedcomment}
A few parameters used in the handover algorithm aren't listed in the main parameter table:
$N_{\rm iter}$ (the maximum number of iterations in the fixed-point refinement step), $\zeta$
(the factor used to relax the capacity constraint when the relaxation is infeasible), and
$\epsilon$ (the convergence tolerance for that iteration) in Algorithm 1. Worth adding for
reproducibility.
\end{quotedcomment}
\begin{authorresponse}
We agree and have added all three parameters to the main simulation table:
$N_{\rm iter}=3$, $\zeta=1.1$, and $\epsilon=1$ RB. In addition, the algorithm section now
explains that $N_{\rm iter}$ caps the frame-level fixed-point computation, $\zeta$ applies
10\% capacity-relaxation increments when the relaxed problem is infeasible, and $\epsilon$
matches the one-RB granularity of the final integer allocation.
\end{authorresponse}
\begin{manuscriptchanges}
Table II now contains the rows ``GAP-HO iteration cap $N_{\rm iter}$,''
``Capacity-relaxation factor $\zeta$,'' and ``Convergence tolerance $\epsilon$,'' with values
$3$, $1.1$, and $1$ RB, respectively.
\end{manuscriptchanges}
\revisionlocation{Page 8, Section IV-B; Page 11, Table II.}

\commentheading{Related work on joint optimization}
\begin{quotedcomment}
The related work section covers handover, beamforming, and resource allocation as separate
lines of work, even though the paper's central claim rests on coordinating them. It would
strengthen the positioning of the contribution to also discuss prior work that combines two or
more of these problems.
\end{quotedcomment}
\begin{authorresponse}
We agree and have supplemented the three separate literature streams with a paragraph organized
around pairwise coordination. The revised discussion contrasts (i) AEPHORA's joint PHO and RA,
(ii) learning-assisted joint scheduling and BF for mmWave vehicular links, and (iii) vertical-HO
association in heterogeneous macro and mmWave networks. It then identifies what remains absent from these designs:
closed-loop three-way coordination of cross-tier association, beam alignment, and queue-aware RA
over the two decision timescales.
\end{authorresponse}
\begin{manuscriptchanges}
We added:

\emph{``Some studies jointly optimize subsets of these decisions. Our conference work AEPHORA
coordinates PHO and RA using predicted mobility, but does not model BF or its training overhead.
Learning-assisted scheduling and BF have also been jointly studied for mmWave vehicular links,
but without cross-tier PHO, while vertical-HO studies in macro/mmWave HetNets focus on association
rather than the closed-loop coordination of association, beam alignment, and queue-aware RA.
Thus, existing pairwise designs do not capture the three-way, two-timescale coupling considered
here.''}
\end{manuscriptchanges}
\revisionlocation{Page 2, Section I (Introduction), related-work discussion.}

\commentheading{Writing quality}
\begin{quotedcomment}
Overall the writing is clear and well organized, no major language concerns.
\end{quotedcomment}
\begin{authorresponse}
We thank the Reviewer for this positive assessment. We have preserved the manuscript's overall
organization while editing the revised and surrounding passages for terminology, consistency,
and readability.
\end{authorresponse}
\revisionlocation{No specific manuscript change is required for this comment.}

\vfill
\begin{flushleft}
Yours sincerely,\\
Bowen Xie, Sheng Zhou, and Zhisheng Niu
\end{flushleft}

\end{document}
